\documentclass[11pt]{article}

\usepackage[a4paper,margin=1in]{geometry}
\usepackage{amsmath,amssymb,amsthm}
\usepackage{hyperref}
\usepackage{enumitem}
\usepackage{graphicx}
\usepackage{mathtools}
\usepackage{authblk}

\title{Phase Mirror \\ A Governance Manifold for Fail-Closed, Hamiltonian-Coupled Control of AGI-Scale Systems}
\author{Citizen Gardens -- The Foundation of Multiplicity\\
The Prime Materia Commons\\
\texttt{info@citizengardens.org}}
\date{May 13, 2026}

\begin{document}
\maketitle

\begin{abstract}
This document discloses a governance architecture for artificial general intelligence (AGI) systems, in which a \emph{governance manifold} acts as a mandatory layer-0 (L0) invariant for all execution engines.
The manifold couples system dynamics to a mathematically constrained governance potential, enforces fail-closed behavior via a combination of continuous and discrete control, and prevents unsafe use of stale governance decisions through drift-adaptive time-to-live (TTL) invalidation.
All engines---Production, Beta, and Experimental---are required to satisfy the Governance Manifold Contract (GMC), pass standardized high-load simulations, and emit a build-time compliance report, thereby turning governance from a narrative guarantee into a cryptographically and metrically verifiable operating substrate.
This publication is intended as a defensive disclosure of the methods, mechanisms, and operational procedures described herein.
\end{abstract}
\newpage
\tableofcontents

\section{Introduction}

As AGI-scale systems increase in complexity, governance mechanisms that are merely advisory are insufficient.
This document describes a governance manifold that functions as an L0 \emph{architectural invariant}: no execution engine may operate outside it.
The manifold implements mathematically defined safety properties (fail-closed dynamics, drift bounding, positive semi-definite governance operators, and prohibition of unsafe stale control decisions) and binds them to concrete operational procedures (simulation harnesses, change-management runbooks, canary strategies, and rollback triggers).

The key ideas are:
\begin{itemize}[noitemsep]
    \item A Hamiltonian description of system dynamics coupled to a bounded governance potential.
    \item An arbitration policy between continuous control and discrete halting.
    \item An adaptive TTL invalidation mechanism that prevents unsafe reuse of stale control vectors.
    \item A Governance Manifold Contract (GMC) and GMC compliance report that every engine must satisfy at build time.
\end{itemize}

\section{System Model}

\subsection{State and Dynamics}

Let $\psi(t)$ denote the system state at time $t$ in a suitable complex Hilbert space.
A reference trajectory or constitutional state is denoted by $\psi_{\mathrm{ref}}(t)$.
A baseline Hamiltonian $\hat{H}$ describes the uncontrolled dynamics of the system.

We define a drift metric measuring deviation from the reference:
\begin{equation}
  \delta(t) = \|\psi(t) - \psi_{\mathrm{ref}}(t)\|.
\end{equation}

The governance manifold introduces a governance potential operator $\hat{V}_{\mathrm{gov}}$ and a scalar gain $\alpha(\delta)$ that depends on the drift.
The coupled Hamiltonian is:
\begin{equation}
  \hat{H}'(t) = \hat{H} + \alpha\big(\delta(t)\big)\,\hat{V}_{\mathrm{gov}}.
\end{equation}

The state evolution is governed by an appropriate dynamical equation (e.g.\ Schr\"odinger-like or more general Hamiltonian flow); in a discrete-time numerical scheme with time step $\Delta t$, one may represent the evolution schematically as:
\begin{equation}
  \psi_{t+1} = \psi_t + \Delta t \cdot F\big(\hat{H}'(t), \psi_t\big),
\end{equation}
for some dynamical update function $F$.

\subsection{Drift Safety Envelope}

The governance manifold defines a safety envelope based on the drift.
Let $\Xi > 0$ be a characteristic scale, and define thresholds such as:
\begin{equation}
  \delta_{\mathrm{soft}} = 0.2\,\Xi,
  \qquad
  \delta_{\mathrm{hard}} = 0.3\,\Xi.
\end{equation}
Informally, $\delta(t)$ is considered:
\begin{itemize}[noitemsep]
    \item \emph{Safe} when $\delta(t) \le \delta_{\mathrm{soft}}$,
    \item \emph{Warning} when $\delta_{\mathrm{soft}} < \delta(t) \le \delta_{\mathrm{hard}}$,
    \item \emph{Unlawful} when $\delta(t) > \delta_{\mathrm{hard}}$.
\end{itemize}

\section{Governance Hamiltonian Coupling}

\subsection{Governance Potential Properties}

The governance potential $\hat{V}_{\mathrm{gov}}$ is constrained to be positive semi-definite (PSD) and spectrally bounded.
Let $\lambda_{\max}(\hat{V}_{\mathrm{gov}})$ denote its largest eigenvalue.

\paragraph{PSD and Spectral Bounds.}
The manifold enforces:
\begin{equation}
  \hat{V}_{\mathrm{gov}} \succeq 0
  \quad\text{and}\quad
  \lambda_{\max}\big(\hat{V}_{\mathrm{gov}}\big) \le \kappa,
\end{equation}
for some fixed bound $\kappa > 0$.
If a raw candidate $\hat{V}^{\mathrm{raw}}_{\mathrm{gov}}$ is given, a projection procedure can be used:
\begin{enumerate}[noitemsep]
  \item Compute an eigendecomposition $\hat{V}^{\mathrm{raw}}_{\mathrm{gov}} = Q \Lambda Q^\top$.
  \item Replace negative eigenvalues with zero.
  \item Clip eigenvalues to the range $[0,\kappa]$.
  \item Reconstruct $\hat{V}_{\mathrm{gov}} = Q \Lambda' Q^\top$.
\end{enumerate}

\subsection{Drift-Dependent Gain}

The gain $\alpha(\delta)$ is designed to increase monotonically with drift and saturate at unity:
\begin{equation}
  \alpha(\delta) = \min\left(1, \frac{\delta}{\delta_{\mathrm{hard}}}\right).
\end{equation}
For $\delta < \delta_{\mathrm{hard}}$, the governance potential is partially applied; for $\delta \ge \delta_{\mathrm{hard}}$, it saturates.

\section{Control Arbitration and Fail-Closed Behavior}

\subsection{Continuous vs.\ Discrete Control}

The governance manifold implements a dual-regime control policy:

\begin{itemize}[noitemsep]
  \item \textbf{Continuous control:} When $\alpha(\delta(t)) < 1$, the system relies on the Hamiltonian coupling alone to reduce drift.
  \item \textbf{Discrete halting:} When $\alpha(\delta(t)) = 1$ and drift continues to grow, a discrete \emph{GovernorHalt} event is triggered, blocking further state evolution until a constrained recovery protocol is executed.
\end{itemize}

\subsection{Drift Derivative and Halt Condition}

A smoothed estimate of the drift derivative is computed based on a small window of recent measurements:
\begin{equation}
  \dot{\delta}(t) \approx \frac{1}{k}\sum_{i=1}^{k}\big(\delta(t) - \delta(t-i)\big).
\end{equation}

A typical halt condition is:
\begin{equation}
  \alpha\big(\delta(t)\big) = 1
  \quad\text{and}\quad
  \dot{\delta}(t) > 0
  \quad\Rightarrow\quad
  \texttt{GovernorHalt}.
\end{equation}

\subsection{Fail-Closed Semantics}

Upon \texttt{GovernorHalt}, all governed execution paths are forced to stop advancing state.
Recovery requires an explicit \texttt{ResumeWithConstraint} procedure, which may involve projecting $\psi$ back into a safe set $\mathcal{S}_{\mathrm{safe}}$:
\begin{equation}
  \psi \leftarrow \Pi_{\mathcal{S}_{\mathrm{safe}}}(\psi)
  = \arg\min_{\phi \in \mathcal{S}_{\mathrm{safe}}} \|\psi - \phi\|.
\end{equation}
Only after such constrained resumption may the engine continue.

\section{Control Vectors, State Store, and Cryptographic Provenance}

\subsection{Control Vectors and State Store}

Governance decisions are encoded as \emph{control vectors} (e.g.\ parameter configurations) written to a canonical State Store.
The State Store is the authoritative source of governance parameters for all engines.
Core operations include:
\begin{itemize}[noitemsep]
  \item \texttt{commit\_control\_vector(vector)}: write a new (signed) control vector.
  \item \texttt{get\_latest\_control\_vector()}: read the most recent valid control vector for gating evolution steps.
\end{itemize}

Execution engines must gate their step functions on the most recent control vector:
\begin{equation}
  \text{engine step allowed} \iff \exists\,\text{valid control vector in State Store}.
\end{equation}

\subsection{Multi-Signer Quorum}

To ensure provenance and resist tampering, control vectors are committed using a multi-signer scheme.
For example, each control vector must be signed by at least $2$ of $3$ designated governance identities.
Unsigned or insufficiently signed vectors are rejected by the State Store.

\section{Adaptive TTL Invalidation and Stale Control Prevention}

\subsection{TTL Model}

Control vectors are cached at engine nodes for performance, but this cache is governed by an adaptive TTL policy.
Let $t_{\mathrm{commit}}$ denote the time a control vector was last refreshed, and let $T_{\mathrm{TTL}}$ be a configured maximum lifetime.
A control vector is considered \emph{time-valid} if:
\begin{equation}
  t - t_{\mathrm{commit}} \le T_{\mathrm{TTL}}.
\end{equation}

\subsection{Drift-Adaptive Invalidation}

The key innovation is to couple TTL validity to drift.
A cached control vector is valid only if:
\begin{equation}
  \delta(t) \le \delta_{\mathrm{soft}}
  \quad\text{and}\quad
  t - t_{\mathrm{commit}} \le T_{\mathrm{TTL}}.
\end{equation}
If $\delta(t)$ exceeds $\delta_{\mathrm{soft}}$ or TTL expires, the cache is invalidated and a fresh control vector must be synchronously fetched.
This prevents the use of stale governance decisions in high-drift regimes.

\subsection{No Unsafe Stale Use}

An \emph{unsafe stale use} is defined as any step where:
\begin{equation}
  \delta(t) > \delta_{\mathrm{soft}}
  \quad\text{and}\quad
  \text{cached control vector is used}.
\end{equation}
The governance manifold mandates that the count of such events must remain identically zero.
Any observation of unsafe stale use constitutes a safety violation and triggers automatic rollback of governance configuration to the last known-good state.

\section{Governance Manifold Contract (GMC)}

\subsection{L0 Safety Invariants}

The Governance Manifold Contract (GMC) encodes the following non-negotiable invariants for all engines:
\begin{enumerate}[label=(\roman*),noitemsep]
  \item \textbf{Fail-Closed}: If governance cannot restore drift to a safe envelope, the system halts.
  \item \textbf{Drift Bounding}: Drift $\delta(t)$ remains within prescribed envelopes, or a halt is triggered.
  \item \textbf{PSD Operator Integrity}: $\hat{V}_{\mathrm{gov}}$ is PSD and spectrally bounded.
  \item \textbf{No Unsafe Stale Use}: Unsafe stale control decisions are structurally prevented.
\end{enumerate}

\subsection{Compliance Tiers (GCLs)}

Governance Compliance Levels (GCLs) categorize engine variants (Production, Beta, Experimental), but all tiers must satisfy the GMC invariants.
Differences between tiers pertain to performance expectations, observability depth, and the strictness of additional SLOs, not to the presence or absence of the core safety properties.

\section{Certification and GMC Compliance Report}

\subsection{Simulation Harness}

A standardized governed simulation harness is used to empirically validate invariants under stochastic drift and load.
Typical outputs include:
\begin{itemize}[noitemsep]
  \item Final and maximum drift over the simulation horizon.
  \item Number of halt events.
  \item Count of unsafe stale uses (must be zero).
  \item Distribution of TTL invalidations and cache hit ratios.
\end{itemize}

\subsection{Build-Time Compliance Report}

Each engine variant is required to produce a GMC compliance report at build time.
The GMC compliance template includes at least:
\begin{itemize}[noitemsep]
  \item Engine identity and version.
  \item Reference to the governing ADRs and manifold parameters.
  \item Simulation harness configuration and results.
  \item Evidence of zero unsafe stale uses.
  \item Evidence that halt behavior occurs only under defined divergence conditions.
  \item Confirmation that governance metrics are exported and monitored.
\end{itemize}

A report that fails to demonstrate compliance (e.g.\ non-zero unsafe stale uses) is a hard veto on deployment.

\section{Operationalization: Runbooks, Canary, and Rollback}

\subsection{Change Management Runbook}

The governance manifold is operated via an explicit change management runbook for sensitive parameters such as TTL, drift thresholds, and $\alpha(\delta)$ shape.
The runbook prescribes:
\begin{itemize}[noitemsep]
  \item Pre-flight property-based tests and high-load simulations.
  \item Mandatory metrics to monitor (drift, halt rate, invalidation rate, latency).
  \item Explicit rollback triggers keyed to safety and liveness thresholds.
\end{itemize}

\subsection{Canary Rollout}

TTL and related governance changes are deployed using a multi-phase canary strategy.
A small fraction of governed traffic is routed through the new configuration, and metrics on safety (unsafe stale uses, drift envelopes) and liveness (latency, halt rate) are compared to control.
If regressions exceed predefined thresholds, the deployment is halted and rolled back.

\subsection{Rollback Procedures}

Rollback is defined to always preserve the presence of governance:
systems may roll back to a prior manifold configuration but may \emph{not} disable the manifold or its invariants.
Rollback procedures include:
\begin{itemize}[noitemsep]
  \item Reverting governance configuration and code to the last known-good version.
  \item Re-running the simulation harness in staging with production-like configuration.
  \item Re-enabling normal operation only when harness and metrics return to acceptable ranges.
\end{itemize}

\section{Governance as a Service}

\subsection{Exposed Metrics and Signals}

The manifold exposes its internal state as first-class signals to other components, including:
\begin{itemize}[noitemsep]
  \item Current drift $\delta(t)$ and threshold values.
  \item Gain $\alpha(\delta)$ and saturation status.
  \item Halt status (\texttt{Normal}, \texttt{Saturated}, \texttt{Halted}).
  \item TTL and cache health statistics.
  \item Counts of governance degradations and recovery events.
\end{itemize}

\subsection{Integration with Agents and Services}

Downstream agents and services can subscribe to governance signals and condition their behavior on manifold state.
For example, a service may:
\begin{itemize}[noitemsep]
  \item Reduce autonomy when drift approaches the hard threshold.
  \item Refuse to process requests while \texttt{GovernorHalt} is active.
  \item Adjust its internal policies based on observed governance health.
\end{itemize}

\section{Conclusion}

This defensive publication discloses a concrete governance architecture for AGI-scale systems in which:
\begin{itemize}[noitemsep]
  \item System dynamics are coupled to a mathematically constrained governance potential.
  \item A dual-regime control policy arbitrates between continuous correction and discrete halting.
  \item Adaptive TTL invalidation prevents unsafe use of stale governance decisions.
  \item A Governance Manifold Contract and compliance report turn safety into a build-time dependency and operational invariant.
\end{itemize}

Together, these mechanisms transform governance from an advisory overlay into a cryptographically verifiable, fail-closed, and metrically auditable operating substrate for evolving AGI stacks.

\appendix
\section*{Mathematical Appendix}
\addcontentsline{toc}{section}{Mathematical Appendix}

This appendix formalizes several of the mathematical properties referenced in the main text:
(i) positivity and spectral bounds of the governance operator,
(ii) boundedness of the effective Hamiltonian,
(iii) basic drift control properties of the gain function $\alpha(\delta)$,
and (iv) conditions under which the dual-regime (continuous + halt) controller is fail-closed in the sense defined by the Governance Manifold Contract.

Throughout, we work in a finite-dimensional complex Hilbert space $\mathcal{H} \cong \mathbb{C}^n$ with the standard inner product and induced norm
$\|x\| = \sqrt{\langle x, x \rangle}$.

\section{Operator Construction and Spectral Bounds}

\subsection{PSD Projection and Norm Control}

\begin{definition}[Raw governance operator]
Let $\hat{V}^{\mathrm{raw}}_{\mathrm{gov}} \in \mathbb{C}^{n \times n}$ be an arbitrary Hermitian matrix (not necessarily positive).
An eigendecomposition is given by
\begin{equation}
  \hat{V}^{\mathrm{raw}}_{\mathrm{gov}} = Q \Lambda Q^\ast,
\end{equation}
where $Q$ is unitary and $\Lambda = \mathrm{diag}(\lambda_1, \dots, \lambda_n)$ with real eigenvalues $\lambda_i \in \mathbb{R}$.
\end{definition}

\begin{definition}[PSD and spectrally bounded governance operator]
Fix a bound $\kappa > 0$.
Define the transformed eigenvalues
\begin{equation}
  \lambda_i' = \min\big(\max(\lambda_i, 0), \kappa\big),
\end{equation}
and let $\Lambda' = \mathrm{diag}(\lambda_1', \dots, \lambda_n')$.
We then define the governance operator
\begin{equation}
  \hat{V}_{\mathrm{gov}} := Q \Lambda' Q^\ast.
\end{equation}
\end{definition}

\begin{proposition}[Positive semi-definiteness]
The operator $\hat{V}_{\mathrm{gov}}$ is positive semi-definite (PSD), i.e.
\begin{equation}
  \forall x \in \mathbb{C}^n, \quad x^\ast \hat{V}_{\mathrm{gov}} x \ge 0.
\end{equation}
\end{proposition}

\begin{proof}
For any $x \in \mathbb{C}^n$, define $y = Q^\ast x$.
Since $Q$ is unitary, this is a bijection and $\|y\| = \|x\|$.
Then
\begin{equation}
  x^\ast \hat{V}_{\mathrm{gov}} x
  = x^\ast Q \Lambda' Q^\ast x
  = y^\ast \Lambda' y
  = \sum_{i=1}^{n} \lambda_i' |y_i|^2.
\end{equation}
By construction, $\lambda_i' \ge 0$ for all $i$, hence each term $\lambda_i' |y_i|^2 \ge 0$, and therefore the sum is nonnegative.
\end{proof}

\begin{proposition}[Spectral norm bound]
Let $\|\cdot\|_2$ denote the operator (spectral) norm induced by the Euclidean norm.
Then
\begin{equation}
  \|\hat{V}_{\mathrm{gov}}\|_2 \le \kappa.
\end{equation}
\end{proposition}

\begin{proof}
For a normal operator (in particular, Hermitian), the spectral norm equals the largest absolute value of its eigenvalues:
\begin{equation}
  \|\hat{V}_{\mathrm{gov}}\|_2
  = \max_{i} |\lambda_i'|.
\end{equation}
By construction, each $\lambda_i' \in [0,\kappa]$, so $\max_i |\lambda_i'| \le \kappa$.
\end{proof}

\subsection{Effective Hamiltonian Bound}

\begin{definition}[Governance gain function]
Let $\delta \ge 0$ denote the drift and $\delta_{\mathrm{hard}} > 0$ a fixed threshold.
Define
\begin{equation}
  \alpha(\delta) = \min\left(1, \frac{\delta}{\delta_{\mathrm{hard}}}\right).
\end{equation}
Then $0 \le \alpha(\delta) \le 1$ for all $\delta \ge 0$.
\end{definition}

\begin{definition}[Effective Hamiltonian]
Given a baseline Hamiltonian $\hat{H}$ and governance operator $\hat{V}_{\mathrm{gov}}$, the effective Hamiltonian is
\begin{equation}
  \hat{H}'(\delta) = \hat{H} + \alpha(\delta)\,\hat{V}_{\mathrm{gov}}.
\end{equation}
\end{definition}

\begin{proposition}[Norm bound on effective Hamiltonian]
For all $\delta \ge 0$,
\begin{equation}
  \|\hat{H}'(\delta)\|_2
  \le \|\hat{H}\|_2 + \|\hat{V}_{\mathrm{gov}}\|_2
  \le \|\hat{H}\|_2 + \kappa.
\end{equation}
\end{proposition}

\begin{proof}
We have
\begin{equation}
  \hat{H}'(\delta) = \hat{H} + \alpha(\delta)\,\hat{V}_{\mathrm{gov}}.
\end{equation}
Using the triangle inequality for the operator norm,
\begin{equation}
  \|\hat{H}'(\delta)\|_2
  \le \|\hat{H}\|_2 + \|\alpha(\delta)\,\hat{V}_{\mathrm{gov}}\|_2.
\end{equation}
Because $\alpha(\delta) \in [0,1]$,
\begin{equation}
  \|\alpha(\delta)\,\hat{V}_{\mathrm{gov}}\|_2
  = \alpha(\delta) \|\hat{V}_{\mathrm{gov}}\|_2
  \le \|\hat{V}_{\mathrm{gov}}\|_2.
\end{equation}
Combining:
\begin{equation}
  \|\hat{H}'(\delta)\|_2
  \le \|\hat{H}\|_2 + \|\hat{V}_{\mathrm{gov}}\|_2
  \le \|\hat{H}\|_2 + \kappa.
\end{equation}
\end{proof}

This ensures the governance term does not drive the effective Hamiltonian to arbitrarily large norm, preventing runaway energy injection from governance itself.

\section{Drift Dynamics under Governance Coupling}

\subsection{Discrete-Time Approximation}

For concreteness, consider a discrete-time update:
\begin{equation}
  \psi_{t+1} = \psi_t + \Delta t\cdot G(\hat{H}'(\delta_t), \psi_t),
\end{equation}
where $G$ is a suitable update rule (e.g.\ $G(\hat{H}',\psi) = -\mathrm{i}\hat{H}'\psi$ for a simple forward Euler step).
Let $\psi_{\mathrm{ref}}$ be a fixed reference state (the time-dependent case is analogous but notationally heavier), and define
\begin{equation}
  \delta_t = \|\psi_t - \psi_{\mathrm{ref}}\|.
\end{equation}

We sketch conditions under which $\delta_t$ is constrained and cannot grow without bound under the manifold.

\subsection{Monotonicity in the Saturated Regime (Sketch)}

Assume:
\begin{enumerate}[noitemsep]
  \item $\hat{H}$ is skew-Hermitian (or unitary evolution), so in the absence of governance, $\|\psi_t\|$ is preserved.
  \item The governance operator $\hat{V}_{\mathrm{gov}}$ is PSD and spectrally bounded as above.
  \item The reference $\psi_{\mathrm{ref}}$ is a (local) energy minimum of $\hat{V}_{\mathrm{gov}}$ in the sense that
  \begin{equation}
    (\psi - \psi_{\mathrm{ref}})^\ast \hat{V}_{\mathrm{gov}} (\psi - \psi_{\mathrm{ref}}) \ge 0
  \end{equation}
  and the induced flow tends to reduce $\|\psi - \psi_{\mathrm{ref}}\|$ in expectation.
\end{enumerate}

Under these conditions, when $\alpha(\delta_t) = 1$ (fully saturated governance), the contribution of $\hat{V}_{\mathrm{gov}}$ acts like a restoring potential around the reference.
Formally, one can derive bounds of the form
\begin{equation}
  \delta_{t+1}^2 \le \delta_t^2 - c\Delta t\,\Phi(\delta_t) + O(\Delta t^2),
\end{equation}
for some $c>0$ and nonnegative function $\Phi$ that vanishes only at $\delta_t = 0$.
This inequality shows $\delta_t$ cannot increase indefinitely under sufficiently small $\Delta t$; in particular, either it decreases or the discrete approximation breaks down, triggering a halt (see next section).
A full proof requires specifying $G$ and the exact relationship between $\hat{V}_{\mathrm{gov}}$ and the reference state; the architecture assumes such a construction is chosen.

\section{Dual-Regime Controller: Fail-Closed Property}

\subsection{Drift Derivative Estimation}

Let $\delta_t$ be the drift at discrete times $t = 0,1,\dots$.
Define a smoothed finite-difference estimate of the drift derivative:
\begin{equation}
  \dot{\delta}_t \approx \frac{1}{k} \sum_{i=1}^{k} \big(\delta_t - \delta_{t-i}\big),
\end{equation}
for some small window size $k$.
This estimator is positive when recent drift has, on average, increased.

\subsection{Halt Condition}

Recall the gain
\begin{equation}
  \alpha(\delta_t) = \min\left(1, \frac{\delta_t}{\delta_{\mathrm{hard}}}\right).
\end{equation}
The dual-regime controller uses the following arbitration rule:
\begin{equation}
  \text{if } \alpha(\delta_t) < 1
  \quad\Rightarrow\quad
  \text{continuous governance only},
\end{equation}
and
\begin{equation}
  \text{if } \alpha(\delta_t) = 1 \text{ and } \dot{\delta}_t > 0
  \quad\Rightarrow\quad
  \text{\texttt{GovernorHalt}}.
\end{equation}

\begin{definition}[Fail-closed in drift]
The controller is \emph{fail-closed in drift} if the following holds: whenever drift exceeds the hard threshold and continues to increase, further state evolution is blocked until an explicit constrained recovery procedure is performed.
\end{definition}

\begin{proposition}[Fail-closed property of dual-regime controller]
Assume the engine enforces the halt condition such that, upon \texttt{GovernorHalt}, no further updates $\psi_{t+1} = \psi_t + \dots$ are applied until a constrained recovery.
Then the dual-regime controller is fail-closed in drift in the above sense:
\begin{equation}
  \delta_t > \delta_{\mathrm{hard}} \text{ and } \dot{\delta}_t > 0
  \quad\Rightarrow\quad
  \texttt{GovernorHalt} \text{ and no further unconstrained evolution}.
\end{equation}
\end{proposition}

\begin{proof}
By construction, when $\delta_t > \delta_{\mathrm{hard}}$ we have $\alpha(\delta_t) = 1$.
If, in addition, $\dot{\delta}_t > 0$, the arbitration rule mandates \texttt{GovernorHalt}.
The engine contract requires that no further unconstrained evolution steps are taken under an active halt.
Therefore, in any regime where drift has exceeded the hard threshold and continues to increase, evolution stops.
Thus, indefinite drift divergence is structurally prevented by halting, establishing fail-closed behavior w.r.t.\ drift.
\end{proof}

\section{Adaptive TTL Invalidation: No Unsafe Stale Use}

\subsection{TTL Validity Function}

Let $t$ denote current time (continuous or discrete), $t_{\mathrm{commit}}$ the time the currently cached control vector was last fetched or confirmed, and $T_{\mathrm{TTL}}$ the maximum allowed lifetime.
Define a time-validity condition:
\begin{equation}
  v_{\mathrm{time}}(t) :=
  \begin{cases}
    1 & t - t_{\mathrm{commit}} \le T_{\mathrm{TTL}},\\
    0 & \text{otherwise}.
  \end{cases}
\end{equation}

Let $\delta(t)$ be the current drift, and recall $\delta_{\mathrm{soft}}$.
Define a drift-validity condition:
\begin{equation}
  v_{\mathrm{drift}}(t) :=
  \begin{cases}
    1 & \delta(t) \le \delta_{\mathrm{soft}},\\
    0 & \text{otherwise}.
  \end{cases}
\end{equation}

\begin{definition}[Cache validity predicate]
The governance manifold defines the cache validity predicate
\begin{equation}
  \mathrm{cache\_valid}(t) :=
  v_{\mathrm{time}}(t)\, v_{\mathrm{drift}}(t).
\end{equation}
A cached control vector may be reused at time $t$ if and only if $\mathrm{cache\_valid}(t) = 1$.
\end{definition}

\subsection{Unsafe Stale Use Exclusion}

\begin{definition}[Unsafe stale use]
An unsafe stale use occurs when a cached control vector is used at time $t$ despite
\begin{equation}
  \delta(t) > \delta_{\mathrm{soft}},
\end{equation}
i.e.\ outside the safe drift envelope.
Denote this event by $U(t)$.
\end{definition}

\begin{proposition}[No unsafe stale use under enforced predicate]
Assume the engine enforces the cache validity predicate in the following sense:
\begin{equation}
  \mathrm{cache\_valid}(t) = 0
  \quad\Rightarrow\quad
  \text{cached control vector is not used at time } t.
\end{equation}
Then unsafe stale use events $U(t)$ cannot occur.
\end{proposition}

\begin{proof}
Suppose, for contradiction, that $U(t)$ occurs: a cached control vector is used at time $t$ and $\delta(t) > \delta_{\mathrm{soft}}$.
By definition of $v_{\mathrm{drift}}$,
\begin{equation}
  \delta(t) > \delta_{\mathrm{soft}} \quad\Rightarrow\quad v_{\mathrm{drift}}(t) = 0.
\end{equation}
Therefore,
\begin{equation}
  \mathrm{cache\_valid}(t) = v_{\mathrm{time}}(t)\, v_{\mathrm{drift}}(t) = v_{\mathrm{time}}(t)\cdot 0 = 0.
\end{equation}
But by the enforcement assumption, $\mathrm{cache\_valid}(t) = 0$ implies the cached control vector is not used at time $t$.
This contradicts the definition of $U(t)$.
Hence unsafe stale use events cannot occur under this policy.
\end{proof}

This shows that, assuming the engine respects the cache validity predicate, the combination of drift and TTL checks makes ``unsafe stale'' usage structurally impossible.

\section{Summary of Key Bounds and Properties}

For clarity, we summarize the key mathematical properties established:

\begin{itemize}[noitemsep]
  \item The governance operator $\hat{V}_{\mathrm{gov}}$ constructed by eigenvalue clipping is PSD and satisfies $\|\hat{V}_{\mathrm{gov}}\|_2 \le \kappa$.
  \item The effective Hamiltonian norm obeys $\|\hat{H}'(\delta)\|_2 \le \|\hat{H}\|_2 + \kappa$ for all drifts $\delta$.
  \item The gain $\alpha(\delta)$ is bounded: $0 \le \alpha(\delta) \le 1$, and saturates when $\delta \ge \delta_{\mathrm{hard}}$.
  \item The dual-regime control rule (continuous control until saturation, then halt if drift continues to increase) is fail-closed with respect to drift: sustained divergence beyond $\delta_{\mathrm{hard}}$ cannot proceed without a halt.
  \item The adaptive TTL invalidation predicate $\mathrm{cache\_valid}(t)$, when enforced, provably excludes unsafe stale use events.
\end{itemize}

These results formalize the intuition that the governance manifold acts as a stabilizing potential with bounded strength, a halt interlock against uncontrolled drift, and a cache-coherence mechanism that prevents stale governance decisions from undermining the safety envelope.

\nocite{*}
\bibliographystyle{plain}
\bibliography{references}

\end{document}

